\begin{proof}[Proof of \cref{obj:thm:paired-uniform-tv-frontier}]
Throughout, resources are admissible:
\[
n,d\in\mathbb N,\qquad n,d\ge2,\qquad 0<\varepsilon\le1.
\]
We use the effective resource, logarithmic dimension, and converse tuning quantities
\[
t=n\varepsilon^2>0,\qquad L=\log(ed),\qquad
z_*=\frac{d^2L}{t},\qquad w=\log(e+z_*).
\]
In particular, \(1\le L\le d\). The rates in \cref{obj:def:frontier-handle} are independent of the protocol-class label:
\[
r_{\mathrm{NI}}=r_{\mathrm{SI}}=\rho(t,d),
\qquad
h_{\mathrm{NI}}=h_{\mathrm{SI}}=\sqrt{\rho(t,d)}.
\]

\begin{enumerate}
\item \textit{A uniform approximation constant.}
Invoke the published Bernstein limit
\[
K e_K\longrightarrow\beta_*
\quad\text{along positive even integers }K,
\qquad 0.278<\beta_*<0.286,
\]
as stated in \citep[Section 3.1, definition of the Bernstein constant and the immediately following paragraph, p.~8]{CaiLow2011Nonsmooth}; see also \citep{Bernstein1913}.
Choose an integer \(N_{\mathrm B}\ge1\) such that
\[
K e_K\ge\frac1{10}
\qquad\text{for every even }K\ge2N_{\mathrm B},
\]
and define
\[
b_{\mathrm{app}}
=\min\left\{\frac1{10},e_{2N_{\mathrm B}}\right\}>0.
\]
The positivity follows from the preceding tail bound at \(K=2N_{\mathrm B}\).
Approximation error decreases as the allowed degree increases. Hence, for every positive even \(K\),
\[
\frac{b_{\mathrm{app}}}{K}\le e_K.
\]
Indeed, the tail bound proves this when \(K\ge2N_{\mathrm B}\); otherwise,
\[
b_{\mathrm{app}}\le e_{2N_{\mathrm B}}\le e_K\le K e_K.
\]
% lean: CausalSmith.Stat.LdpOptvalueUniformFrontier.bernstein_uniform_positive_of_gate

\item \textit{Paired product priors.}
Take the converse amplitude and matching degree from \cref{obj:def:frontier-handle}:
\[
(\alpha,k)=
\begin{cases}
\left(\dfrac{\sqrt{z_*}}{100},\,2\lceil16L\rceil+1\right),&z_*\le1,\\[4pt]
\left(\dfrac1{100},\,2\lceil16L/w\rceil+1\right),&z_*>1,
\end{cases}
\qquad K=k-1.
\]
These choices satisfy \(0<\alpha\le1/2\), and \(K\) is a positive even integer.
By \cref{obj:lem:cai-low-moment-duality}, there are symmetric probability measures \(\xi_0,\xi_1\), supported on \([-1,1]\), whose moments agree through degree \(K\) and whose absolute moments differ by \(2e_K\). Define the scaled coordinate priors and their product priors by
\[
\nu_\ell=(x\mapsto\alpha x)_\#\xi_\ell,
\qquad
\Pi_\ell=\nu_\ell^{\otimes d},
\qquad \ell\in\{0,1\}.
\]
They are supported on \([-\alpha,\alpha]^d\subseteq[-1/2,1/2]^d\), and
\[
\int u^q\,d\nu_1(u)=\int u^q\,d\nu_0(u)
\qquad(0\le q\le K).
\]
Using the paired-target identity in \cref{obj:prop:signed-causal-bridge}, define their total-variation centers and separation by
\[
v_\ell^{\mathrm{TV}}=\frac12\int |u|\,d\nu_\ell(u),
\qquad
\Delta_{\mathrm{TV}}
=v_1^{\mathrm{TV}}-v_0^{\mathrm{TV}}
=\alpha e_K>0.
\]
Independence of the coordinates gives
\[
\mathbb E_{\Pi_\ell}D_{\mathrm{TV}}(P_\theta)=v_\ell^{\mathrm{TV}},
\qquad
\operatorname{Var}_{\Pi_\ell}\!\left(D_{\mathrm{TV}}(P_\theta)\right)
\le\frac{\alpha^2}{4d}.
\]
Consequently, Chebyshev's inequality and the approximation bound give
\[
\Pi_\ell\left\{
\left|D_{\mathrm{TV}}(P_\theta)-v_\ell^{\mathrm{TV}}\right|
\ge\frac{\Delta_{\mathrm{TV}}}{8}
\right\}
\le
\frac{16\alpha^2}{d\Delta_{\mathrm{TV}}^2}
\le\frac{16K^2}{b_{\mathrm{app}}^2d}.
\]

The ceiling inequalities also give the required rate calibration. If \(z_*\le1\), then
\[
K\le34L,\qquad
\sqrt{\rho(t,d)}=\frac{\sqrt{z_*}}{L},
\qquad
\frac{\alpha}{K}\ge\frac1{3400}\sqrt{\rho(t,d)}.
\]
If \(z_*>1\), then
\[
K\le34\max\{1,L/w\},
\qquad
\sqrt{\rho(t,d)}=\min\{1,w/L\},
\]
so again
\[
\frac{\alpha}{K}\ge\frac1{3400}\sqrt{\rho(t,d)}.
\]
Since \(w\ge1\), both branches satisfy \(K\le34L\). Define
\[
s_{\mathrm{app}}=\frac{b_{\mathrm{app}}}{3400}>0.
\]
Then
\[
\Delta_{\mathrm{TV}}\ge s_{\mathrm{app}}\sqrt{\rho(t,d)}.
\]
Finally, \(L^2/d\to0\) as \(d\to\infty\). Thus there is a universal integer \(D_{\mathrm{cut}}\ge2\) such that
\[
\frac{16\cdot34^2L^2}{b_{\mathrm{app}}^2d}\le\frac1{100}
\qquad(d>D_{\mathrm{cut}}).
\]
For those dimensions, each prior's target-concentration failure probability is at most \(1/100\).
% lean: CausalSmith.Stat.LdpOptvalueUniformFrontier.paired_lawspace_resource_priors_of_gate

\item \textit{Contraction of the paired transcript mixtures.}
Fix any sequential paired-input protocol \(\mathsf K\).
By \cref{obj:prop:signed-causal-bridge}, its pullback
\[
Q_i(E\mid(j,a,y),\eta,r)
=
\mathsf K_i\!\left(E\mid(j,(2a-1)(2y-1)),\eta,r\right)
\]
is sequentially private at the same budget, and its seed-transcript law under \(G_\theta\) equals that of \(\mathsf K\) under \(P_\theta\).
The density hypothesis for \cref{obj:lem:adaptive-moment-contraction} is supplied by \cref{obj:lem:measurable-kernel-radon-nikodym}.
Define the paired seed-transcript mixtures by
\[
\mathsf M_\ell^{\mathrm{TV}}
=
\int\mathcal L_{P_\theta,\mathsf K}(R,\mathbf Z^{\mathsf K})\,
d\Pi_\ell(\theta),
\qquad \ell\in\{0,1\},
\]
and the contraction scale by
\[
\beta=\frac{e^\varepsilon-1}{2d}.
\]
The moment agreement in the previous step allows \cref{obj:lem:adaptive-moment-contraction} to yield
\[
\operatorname{TV}(\mathsf M_0^{\mathrm{TV}},\mathsf M_1^{\mathrm{TV}})
\le
\min\left\{1,d\alpha^k\sqrt{\binom nk}\,\beta^k\right\}
\qquad(k\le n),
\]
whereas the two mixtures are equal when \(k>n\).

For \(0<\varepsilon\le1\), the exponential bound
\(e^\varepsilon-1\le2\varepsilon\) gives \(\beta\le\varepsilon/d\).
Moreover,
\[
\binom nk\le\frac{n^k}{k!}\le\left(\frac{en}{k}\right)^k,
\]
because \(k^k/k!\le e^k\), as follows from the exponential series.
Thus, when \(k\le n\),
\[
\operatorname{TV}(\mathsf M_0^{\mathrm{TV}},\mathsf M_1^{\mathrm{TV}})
\le
\min\left\{1,
d\left(\alpha\sqrt{\frac{en}{k}}\frac{\varepsilon}{d}\right)^k
\right\}.
\]

To evaluate this bound, define
\[
w_{\mathrm{con}}=
\begin{cases}
1,&z_*\le1,\\
w,&z_*>1,
\end{cases}
\qquad
B_{\mathrm{con}}=\alpha\sqrt{\frac{en}{k}}\frac{\varepsilon}{d}.
\]
The prescribed ceilings give
\[
32L\le k w_{\mathrm{con}},
\qquad
B_{\mathrm{con}}^2=\frac{\alpha^2eL}{kz_*}.
\]
In the branch \(z_*\le1\), the amplitude choice and \(e<3\) imply
\[
\alpha^2e\,w_{\mathrm{con}}e^{w_{\mathrm{con}}/2}
=\frac{z_*}{10^4}e^{3/2}\le32z_*.
\]
In the other branch, \(e^w=e+z_*\), and \(w\le2e^{w/2}\) gives
\[
w e^{w/2}\le2(e+z_*)\le8z_*.
\]
Hence the same bound
\[
\alpha^2e\,w_{\mathrm{con}}e^{w_{\mathrm{con}}/2}\le32z_*
\]
holds. Combining it with \(32L\le kw_{\mathrm{con}}\) yields
\[
B_{\mathrm{con}}^2\le e^{-w_{\mathrm{con}}/2},
\qquad
B_{\mathrm{con}}\le e^{-w_{\mathrm{con}}/4}.
\]
Therefore
\[
d B_{\mathrm{con}}^k
\le d e^{-8L}
\le e^{-7L}
\le e^{-7}\le\frac1{100}.
\]
For \(k>n\), exact equality of the mixtures gives the same conclusion:
\[
\operatorname{TV}(\mathsf M_0^{\mathrm{TV}},\mathsf M_1^{\mathrm{TV}})
\le\frac1{100}.
\]
% lean: CausalSmith.Stat.LdpOptvalueUniformFrontier.paired_lawspace_resource_contraction

\item \textit{Decision lower bounds in large dimension.}
We give the paired-model decision reduction explicitly.
For an estimator and a globally \(0.90\)-honest interval, define their worst-law risk and expected length by
\[
\begin{aligned}
\mathcal R_{\mathsf K}(\widehat D)
&=\sup_{\mathsf p\in\mathcal P_d^{\mathrm{pair}}}
\mathbb E_{\mathsf p,\mathsf K}
\left[(\widehat D-D_{\mathrm{TV}}(\mathsf p))^2\right],\\
\mathcal H_{\mathsf K}(J_{\mathrm{TV}})
&=\sup_{\mathsf p\in\mathcal P_d^{\mathrm{pair}}}
\mathbb E_{\mathsf p,\mathsf K}\operatorname{length}J_{\mathrm{TV}}.
\end{aligned}
\]
These definitions apply to every admissible Borel decision of \cref{obj:def:tv-private-decision-values}.
For this reduction, let \(\eta_0,\omega\ge0\) bound the two concentration failures and the seed-transcript total variation, respectively.
Let
\[
N_\ell^{\mathrm{TV}}
=\mathsf M_\ell^{\mathrm{TV}}\otimes\operatorname{Uniform}[0,1]
\]
be the joint seed-transcript-randomizer mixture.
Appending the independent analyst randomizer preserves the bound:
\[
\operatorname{TV}(N_0^{\mathrm{TV}},N_1^{\mathrm{TV}})\le\omega.
\]
Indeed, integrating an event indicator over the independent randomizer produces a measurable function taking values in \([0,1]\).

Threshold \(\widehat D\) at
\((v_0^{\mathrm{TV}}+v_1^{\mathrm{TV}})/2\).
On the target-concentration event, a wrong prior label entails an estimation error of magnitude at least \(3\Delta_{\mathrm{TV}}/8\).
If the worst-law risk is finite, Markov's inequality bounds the sum of testing errors by
\[
2\eta_0+
\frac{128\,\mathcal R_{\mathsf K}(\widehat D)}
{9\Delta_{\mathrm{TV}}^2}.
\]
The testing lower bound from total variation is \(1-\omega\). It follows that
\[
\mathcal R_{\mathsf K}(\widehat D)
\ge
\frac{9\Delta_{\mathrm{TV}}^2}{128}
(1-\omega-2\eta_0)_+.
\]
The same inequality holds immediately when the risk is infinite.

For the interval reduction, define
\[
\mathcal E_\ell^{\mathrm{TV}}
=
\left\{
J_{\mathrm{TV}}\cap
\left[v_\ell^{\mathrm{TV}}-\frac{\Delta_{\mathrm{TV}}}{8},
      v_\ell^{\mathrm{TV}}+\frac{\Delta_{\mathrm{TV}}}{8}\right]
\ne\varnothing
\right\}.
\]
Paired-model honesty and target concentration imply
\[
N_\ell^{\mathrm{TV}}(\mathcal E_\ell^{\mathrm{TV}})
\ge0.90-\eta_0.
\]
Transferring the second event to the first mixture and applying the intersection bound gives
\[
N_0^{\mathrm{TV}}
(\mathcal E_0^{\mathrm{TV}}\cap\mathcal E_1^{\mathrm{TV}})
\ge0.80-2\eta_0-\omega.
\]
On this intersection, ordered endpoints and connectedness force length at least \(3\Delta_{\mathrm{TV}}/4\). Therefore
\[
\mathcal H_{\mathsf K}(J_{\mathrm{TV}})
\ge
\frac{3\Delta_{\mathrm{TV}}}{4}
(0.80-2\eta_0-\omega)_+.
\]

For \(d>D_{\mathrm{cut}}\), the preceding steps permit
\(\eta_0=\omega=1/100\). Define
\[
c_{\mathrm{large}}
=\min\left\{\frac{s_{\mathrm{app}}^2}{100},
            \frac{s_{\mathrm{app}}}{100}\right\}>0.
\]
Since \(\Delta_{\mathrm{TV}}\ge s_{\mathrm{app}}\sqrt{\rho(t,d)}\),
\[
\begin{aligned}
\mathcal R_{\mathsf K}(\widehat D)
&\ge\frac9{128}\frac{97}{100}\Delta_{\mathrm{TV}}^2
\ge c_{\mathrm{large}}\rho(t,d),\\
\mathcal H_{\mathsf K}(J_{\mathrm{TV}})
&\ge\frac34\frac{77}{100}\Delta_{\mathrm{TV}}
\ge c_{\mathrm{large}}\sqrt{\rho(t,d)}.
\end{aligned}
\]
% lean: CausalSmith.Stat.LdpOptvalueUniformFrontier.paired_converse_of_gate

\item \textit{Elementary rate comparisons and bounded dimension.}
We first establish the numerical comparisons used to complete the converse:
\[
\min\{1,t^{-1}\}\le\rho(t,d)\le4\min\{1,d^2/t\},
\qquad
\frac14\min\{t,d\}\le t\rho(t,d).
\]
In the dense branch \(t\ge d^2L\), these follow from
\[
\rho(t,d)=\frac{d^2}{tL},\qquad
\frac1t\le\rho(t,d)\le\min\{1,d^2/t\},
\qquad
t\rho(t,d)=\frac{d^2}{L}\ge d.
\]

For the nondense branch, define the auxiliary function
\[
\Psi(s_{\mathrm{res}})
=s_{\mathrm{res}}
\left[\log\left(e+\frac{d^2L}{s_{\mathrm{res}}}\right)\right]^2,
\qquad s_{\mathrm{res}}>0.
\]
Writing \(v_{\mathrm{res}}=d^2L/s_{\mathrm{res}}>0\), differentiation gives
\[
\Psi'(s_{\mathrm{res}})
=
\log(e+v_{\mathrm{res}})
\left[
\log(e+v_{\mathrm{res}})
-\frac{2v_{\mathrm{res}}}{e+v_{\mathrm{res}}}
\right]\ge0.
\]
For the last inequality, \(\log y\ge1-y^{-1}\), applied at
\(y=(e+v_{\mathrm{res}})/e\), gives
\[
\log(e+v_{\mathrm{res}})
\ge1+\frac{v_{\mathrm{res}}}{e+v_{\mathrm{res}}}
\ge\frac{2v_{\mathrm{res}}}{e+v_{\mathrm{res}}}.
\]
Also,
\[
\log(e+d^2L)\ge L,
\qquad
\log(e+dL)\ge L/2.
\]
The first follows from \(dL\ge e\); the second follows by combining
\(\log(e+dL)\ge1\) and \(\log(e+dL)\ge\log d\).
Thus, if \(t\le1\), then \(w\ge L\) and \(\rho(t,d)=1\).
If \(t\ge1\), monotonicity of \(\Psi\) gives
\(tw^2\ge L^2\), establishing the lower comparison with \(\min\{1,t^{-1}\}\).
Likewise, if \(t\le d\), then \(w\ge L/2\); if \(t\ge d\), monotonicity gives
\(tw^2\ge dL^2/4\). Applying the cap at one proves the comparison with \(\min\{t,d\}/4\).

For the upper comparison, \(z_*>1\) in this branch and
\[
w\le\log(e+1)+\log z_*
\le2+2(\sqrt{z_*}-1)=2\sqrt{z_*}.
\]
Hence
\[
\frac{w^2}{L^2}\le\frac{4z_*}{L^2}
=\frac{4d^2}{tL}\le\frac{4d^2}{t}.
\]
Together with \(\rho(t,d)\le1\), this proves all three comparisons.

Now define the two-point amplitude
\[
\alpha_{\mathrm{two}}
=\frac18\min\{1,t^{-1/2}\}\in(0,1/2].
\]
Consider the paired laws with constant contrast zero and constant contrast
\(\alpha_{\mathrm{two}}\). Their targets are zero and
\(\alpha_{\mathrm{two}}/2\).
Applying \cref{obj:lem:adaptive-moment-contraction} at matching order one through the bridge in \cref{obj:prop:signed-causal-bridge} gives
\[
\operatorname{TV}\!\left(
\mathcal L_{P_0,\mathsf K}(R,\mathbf Z^{\mathsf K}),
\mathcal L_{P_{\alpha_{\mathrm{two}}\mathbf1},\mathsf K}
(R,\mathbf Z^{\mathsf K})
\right)
\le
\alpha_{\mathrm{two}}\sqrt n\,\varepsilon
\le\frac18.
\]
Here \(P_{\alpha_{\mathrm{two}}\mathbf1}\) denotes the law with
\(\theta_j=\alpha_{\mathrm{two}}\) for every \(j\).
The preceding decision reduction, with deterministic targets and
\(\omega=1/8\), therefore yields
\[
\mathcal R_{\mathsf K}(\widehat D)
\ge\frac{63\alpha_{\mathrm{two}}^2}{4096}
\ge\frac1{8192}\min\{1,t^{-1}\},
\]
\[
\mathcal H_{\mathsf K}(J_{\mathrm{TV}})
\ge\frac{81\alpha_{\mathrm{two}}}{320}
\ge\frac1{8192}\min\{1,t^{-1/2}\}.
\]
The calibrations use
\[
\alpha_{\mathrm{two}}^2
=\frac1{64}\min\{1,t^{-1}\}.
\]

If \(d\le D_{\mathrm{cut}}\), the elementary upper comparison implies
\[
\rho(t,d)
\le4D_{\mathrm{cut}}^2\min\{1,t^{-1}\},
\]
and taking square roots gives
\[
\sqrt{\rho(t,d)}
\le2D_{\mathrm{cut}}\min\{1,t^{-1/2}\}
\le4D_{\mathrm{cut}}^2\min\{1,t^{-1/2}\}.
\]
Define
\[
c_{\mathrm{small}}=\frac1{32768D_{\mathrm{cut}}^2},
\qquad
c_{\mathrm{lower}}=\min\{c_{\mathrm{large}},c_{\mathrm{small}}\}>0.
\]
The cases \(d\le D_{\mathrm{cut}}\) and \(d>D_{\mathrm{cut}}\) now give, for every sequential paired protocol and every admissible decision,
\[
\mathcal R_{\mathsf K}(\widehat D)\ge
c_{\mathrm{lower}}\rho(t,d),
\qquad
\mathcal H_{\mathsf K}(J_{\mathrm{TV}})\ge
c_{\mathrm{lower}}\sqrt{\rho(t,d)}.
\]
Every noninteractive protocol is sequential. Taking the decision and protocol infima in \cref{obj:def:tv-private-decision-values} therefore proves these lower bounds for both classes.
% lean: CausalSmith.Stat.LdpOptvalueUniformFrontier.paired_minimax_lower_of_gate

\item \textit{The concrete upper-risk construction.}
Choose the universal constants
\[
c=\min\{c_{\mathrm{lower}},1\},
\qquad C_0=10^{16}.
\]
Then \(0<c\le C_0\).
Use the paired procedure specified in \cref{obj:synth_5}, with a singleton public seed. Its channel quantities are
\[
m=\lfloor n/3\rfloor,\qquad
\delta=\tanh(\varepsilon/2),\qquad b=\frac d\delta.
\]
Each of the first \(2m\) participants releases a vector with probabilities
\[
\mathsf K_i(z'\mid(j,s))
=2^{-d}(1+\delta s z'_j),
\qquad z'\in\{-1,1\}^d;
\]
the remaining participants release \((-1,\ldots,-1)\).
Symmetry of the sign-vector alphabet makes the active probabilities sum to one, and
\[
\frac{1+\delta}{1-\delta}=e^\varepsilon
\]
bounds the likelihood ratio at every active message. The constant releases satisfy the privacy inequality as well. All kernels are independent of history, so this is a pure \(\varepsilon\)-private noninteractive protocol with the declared sign-vector message spaces.

For \(n>2\), define the pilot and evaluation matrices, means, and calibration quantities by
\[
W^{\mathrm{pil}}_{ij}=bZ_{i,j},\qquad
W^{\mathrm{ev}}_{ij}=bZ_{m+i,j},
\qquad 1\le i\le m,\quad j\in[d],
\]
\[
M_j=\frac1m\sum_iW^{\mathrm{pil}}_{ij},\qquad
\overline W_j=\frac1m\sum_iW^{\mathrm{ev}}_{ij},
\qquad
\sigma^2=\frac{b^2}{m},\qquad \sigma=\sqrt{\sigma^2},
\qquad H=\log(e+\sigma^2L).
\]
The statistic \(\widehat F\) is the ordered branch-selected norm estimate in
\cref{obj:synth_5}, using the private moments and polynomial coefficients specified there and in \cref{obj:def:private-moments,obj:def:frontier-handle}.
Set
\[
\widehat D=
\begin{cases}
1/8,&n=2,\\
\max\{0,\min\{1/4,\widehat F/2\}\},&n>2.
\end{cases}
\]
The decisions are measurable functions of the finite transcript.

The cube target belongs to \([0,1/4]\), so
\[
(\widehat D-D_{\mathrm{TV}}(P_\theta))^2\le\frac1{16}.
\]
For \(n>2\), clipping also gives the sharper comparison
\[
\mathbb E_{P_\theta,\mathsf K}
(\widehat D-D_{\mathrm{TV}}(P_\theta))^2
\le\frac14
\mathbb E_{P_\theta,\mathsf K}(\widehat F-F(\theta))^2.
\]
To apply the existing moment and polynomial calibrations, use \(G_\theta\) from
\cref{obj:def:symmetric-law}. Its contrasts satisfy
\[
\tau_j(G_\theta)
=\frac{1+\theta_j}{2}-\frac{1-\theta_j}{2}
=\theta_j.
\]
The bridge in \cref{obj:prop:signed-causal-bridge} identifies the active paired blocks with independent vector-channel blocks under this law. Thus
\cref{obj:lem:private-moment-and-dependence,obj:lem:finite-private-polynomial-calibration}
apply with contrast vector \(\theta\); their concentration and approximation hypotheses are supplied by
\cref{obj:lem:bounded-mean-concentration,obj:lem:cai-low-chebyshev-approximation}.

We verify the resource calibration used in the risk bounds. For \(n\ge3\),
\[
m\ge n/5,\qquad
\delta=\frac{e^\varepsilon-1}{e^\varepsilon+1}\ge\varepsilon/4,
\qquad
\sigma^2\le\frac{80d^2}{t}.
\]
The attenuation bound follows from \(e^\varepsilon\ge1+\varepsilon\) and
\(e^\varepsilon+1<4\).
Consequently,
\[
1\le H\le\log(80(e+z_*))
\le5+w\le6w.
\]

We now follow the risk case split.

If \(n=2\), then \(t\le2\), and the elementary lower comparison gives
\(\rho(t,d)\ge1/2\). The constant estimate satisfies
\[
\mathbb E(\widehat D-D_{\mathrm{TV}}(P_\theta))^2
\le\frac1{64}\le C_0\rho(t,d).
\]

Suppose \(n>2\). If \(\rho(t,d)=1\), the uniform \(1/16\) loss bound gives the desired result.
Next consider the constant-fallback branch. Its conditions include
\[
t<d^2L,\qquad \frac{L}{1024H}<4.
\]
Since \(H\le6w\),
\[
L<4096H\le24576w,
\qquad
\rho(t,d)\ge\left(\frac1{24576}\right)^2.
\]
Again,
\[
\mathbb E(\widehat D-D_{\mathrm{TV}}(P_\theta))^2
\le\frac1{16}\le C_0\rho(t,d).
\]

For the remaining branches, suppose first that \(L<4096\).
In the dense case \(t\ge d^2L\), the absolute-mean statistic satisfies
\[
\mathbb E\left[
\left(\frac1d\sum_j|\overline W_j|-F(\theta)\right)^2
\right]
\le\frac1d\sum_j\mathbb E(\overline W_j-\theta_j)^2
\le\sigma^2.
\]
Here the first inequality uses the absolute-value Lipschitz inequality and the quadratic mean inequality; the second uses the moment bounds and independence across participants. The numerical calibration is
\[
\frac{24}{t}+\frac{\sigma^2}{2}
\le\frac{64d^2}{t}
=64L\rho(t,d)
\le64\cdot4096\,\rho(t,d).
\]
Since the TV risk is at most \(\sigma^2/4\), this proves its \(C_0\rho(t,d)\) bound.
In the nondense case, \(w\ge1\) implies
\[
\rho(t,d)\ge\left(\frac1{4096}\right)^2,
\]
so the uniform loss bound suffices.

Finally, suppose \(L\ge4096\).
In the dense case, the prescribed hybrid degree and radius are
\[
D=2\lfloor L/2048\rfloor,\qquad
T_*=4\sigma\sqrt L,\qquad a_{\mathrm{rad}}=8\sigma\sqrt L.
\]
Its degree fits the block:
\[
2D\le L/512\le m.
\]
Indeed, \(n\ge t\ge d^2L\ge4L\) and \(m\ge n/5\).
The hybrid conclusion of \cref{obj:lem:finite-private-polynomial-calibration} gives
\[
\mathbb E(\widehat F-F(\theta))^2
\le500000000\,\frac{\sigma^2}{L}
\le40000000000\,\rho(t,d).
\]
Also \(t^{-1}\le\rho(t,d)\), and hence
\[
\frac{24}{t}
+\frac12\left(500000000\,\frac{\sigma^2}{L}\right)
\le(24+20000000000)\rho(t,d)
\le C_0\rho(t,d).
\]
The clipping comparison proves the claimed TV risk bound.

In the nondense case remaining after the saturation and fallback cases,
\[
\rho(t,d)=(w/L)^2,\qquad
a_{\mathrm{rad}}=\frac12,\qquad
D=2\left\lfloor\frac{L}{2048H}\right\rfloor,
\]
and floor rounding gives
\[
2\le D,\qquad
\frac{L}{2048H}\le D\le\frac{L}{1024H}.
\]
The moment-fit condition \(2D\le m\) also holds. To see this, if \(m<2D\), the degree budget and \(H\ge1\) imply \(m<L/512\).
Since \(\delta<1\),
\[
\sigma^2\ge\frac{d^2}{m},
\qquad
e+\sigma^2L>512d^2\ge ed,
\]
so \(H\ge L\), contradicting \(1024DH\le L\) and \(D\ge2\).

The global-polynomial conclusion of
\cref{obj:lem:finite-private-polynomial-calibration} now yields
\[
\mathbb E(\widehat F-F(\theta))^2
\le\frac{e^{9DH}}{2d}+\frac1{(D+1)^2}.
\]
The degree budget gives \(9DH\le L/100\).
For \(L\ge4096\), the elementary exponential inequalities
\[
L^2\le16e^{L/2},
\qquad
4\le e^{49L/100-1}
\]
imply
\[
L^2e^{L/100}\le4e^{L-1}=4d.
\]
Therefore
\[
\frac{e^{9DH}}{2d}\le\frac2{L^2}\le2\rho(t,d),
\qquad
\frac1{D+1}\le\frac{2048H}{L}
\le\frac{12288w}{L},
\]
and hence
\[
\frac1{(D+1)^2}\le150994944\,\rho(t,d).
\]
The elementary lower comparison, together with \(\rho(t,d)<1\), gives
\(t^{-1}\le\rho(t,d)\). Thus
\[
\frac{24}{t}
+\frac12\left(\frac{e^{9DH}}{2d}+\frac1{(D+1)^2}\right)
\le75497500\,\rho(t,d)
\le C_0\rho(t,d).
\]
Clipping again completes the risk bound. Every admissible tuple has therefore been covered:
\[
\sup_{\mathsf p\in\mathcal P_d^{\mathrm{pair}}}
\mathbb E_{\mathsf p,\mathsf K}
(\widehat D-D_{\mathrm{TV}}(\mathsf p))^2
\le C_0\rho(t,d).
\]
% lean: CausalSmith.Stat.LdpOptvalueUniformFrontier.tvFrontierEstimator_uniform_rate

\item \textit{Intervals and the common minimax bounds.}
Define
\[
q_*=\sqrt{10C_0\rho(t,d)}=\sqrt{10^{17}\rho(t,d)},
\qquad
J_{\mathrm{TV}}
=
\left[
\max\{0,\widehat D-q_*\},
\min\{1/4,\widehat D+q_*\}
\right].
\]
The elementary lower comparison gives \(\rho(t,d)>0\).
The interval contains its center, has ordered Borel endpoints, and is connected, closed, and contained in \([0,1/4]\).
Since the target is also in this range, failure of coverage entails
\(|\widehat D-D_{\mathrm{TV}}(\mathsf p)|>q_*\). Markov's inequality gives
\[
\Pr_{\mathsf p,\mathsf K}
\{D_{\mathrm{TV}}(\mathsf p)\notin J_{\mathrm{TV}}\}
\le\frac{C_0\rho(t,d)}{q_*^2}
=\frac1{10}.
\]
Furthermore,
\[
\operatorname{length}J_{\mathrm{TV}}
\le2q_*
=2\sqrt{10C_0}\sqrt{\rho(t,d)}
\le C_0\sqrt{\rho(t,d)}.
\]
This proves all three concrete attainment guarantees.
The same construction supplies the general attainment witness: its message spaces are exactly \(\{-1,1\}^d\), and its estimate depends on the transcript alone.

The construction belongs to both protocol classes. Its interval, being contained in \([0,1/4]\), is also an admissible interval in the optimization over \([0,1]\).
Taking the corresponding candidates in the minimax infima gives, for either \(\mathcal C\),
\[
\mathfrak R^{\mathrm{TV}}_{\mathcal C}
(\mathcal P_d^{\mathrm{pair}};n,d,\varepsilon)
\le C_0\rho(t,d),
\qquad
\mathfrak H^{\mathrm{TV}}_{\mathcal C}
(\mathcal P_d^{\mathrm{pair}};n,d,\varepsilon)
\le C_0\sqrt{\rho(t,d)}.
\]
Since \(c\le c_{\mathrm{lower}}\) and both rates are nonnegative, the lower bounds proved above give
\[
\begin{aligned}
c\rho(t,d)
&\le\mathfrak R^{\mathrm{TV}}_{\mathcal C}
(\mathcal P_d^{\mathrm{pair}};n,d,\varepsilon)
\le C_0\rho(t,d),\\
c\sqrt{\rho(t,d)}
&\le\mathfrak H^{\mathrm{TV}}_{\mathcal C}
(\mathcal P_d^{\mathrm{pair}};n,d,\varepsilon)
\le C_0\sqrt{\rho(t,d)}.
\end{aligned}
\]
% lean: CausalSmith.Stat.LdpOptvalueUniformFrontier.tvAttainment_minimax_upper

\item \textit{Identification of the paired functional.}
For \(d\ge2\) and \(\theta\in[-1/2,1/2]^d\), the masses
\[
P_\theta(j,s)=\frac{1+s\theta_j}{2d}
\]
are nonnegative, and their sum is one because each pair has total mass \(1/d\).
Taking zero contrast gives the uniform probability law.
For two probability laws on a finite alphabet, total variation equals half the sum of the absolute mass differences. Therefore
\[
D_{\mathrm{TV}}(P_\theta)
=\operatorname{TV}(P_\theta,\mathsf U_{2d}).
\]
For each coordinate,
\[
\sum_{s\in\{-1,1\}}
\left|\frac{1+s\theta_j}{2d}-\frac1{2d}\right|
=
\frac{|\theta_j|}{d}.
\]
Summing and multiplying by one half gives
\[
D_{\mathrm{TV}}(P_\theta)
=\frac1{2d}\sum_{j=1}^d|\theta_j|
=\frac12F(\theta),
\qquad
0\le D_{\mathrm{TV}}(P_\theta)\le\frac14.
\]
% lean: CausalSmith.Stat.LdpOptvalueUniformFrontier.tvFromUniform_pairedLaw

\item \textit{Attainment and protocol-class ratios.}
The witnesses constructed above establish the existential attainment assertions for either class, because the class-selected rates agree. They also establish the concrete guarantees for the procedure of
\cref{obj:synth_5,obj:def:frontier-handle,obj:thm:uniform-private-value-frontiers}.

The protocol inclusion in \cref{obj:def:tv-private-decision-values} gives
\[
\mathfrak R^{\mathrm{TV}}_{\mathrm{SI}}
(\mathcal P_d^{\mathrm{pair}};n,d,\varepsilon)
\le
\mathfrak R^{\mathrm{TV}}_{\mathrm{NI}}
(\mathcal P_d^{\mathrm{pair}};n,d,\varepsilon),
\]
and the corresponding inequality for honest length.
The common bounds make all four values finite and strictly positive.
For risk, they also give
\[
\mathfrak R^{\mathrm{TV}}_{\mathrm{NI}}
(\mathcal P_d^{\mathrm{pair}};n,d,\varepsilon)
\le C_0\rho(t,d)
=\frac{C_0}{c}\,c\rho(t,d)
\le\frac{C_0}{c}
\mathfrak R^{\mathrm{TV}}_{\mathrm{SI}}
(\mathcal P_d^{\mathrm{pair}};n,d,\varepsilon).
\]
Dividing by the positive SI value yields the asserted risk ratio bounds.
Replacing \(\rho(t,d)\) by \(\sqrt{\rho(t,d)}\) proves the same ratio bounds for honest length.
% lean: CausalSmith.Stat.LdpOptvalueUniformFrontier.frontier_ratio_comparison

\item \textit{Saturation and consistency of the rate.}
For \(0<t<d^2L\), positivity of \(w/L\) gives
\[
\begin{aligned}
\rho(t,d)=1
&\Longleftrightarrow w\ge L\\
&\Longleftrightarrow e+\frac{d^2L}{t}\ge ed\\
&\Longleftrightarrow
t\le\frac{d^2L}{ed-e}.
\end{aligned}
\]
The last denominator is positive for every \(d\ge2\).

For an admissible resource sequence, define
\[
t_k=n_k\varepsilon_k^2,\qquad
L_k=\log(ed_k),\qquad
g_k=\frac{\log(1+d_k^2/t_k)}{L_k}.
\]
If \(\rho(t_k,d_k)\to0\), the lower comparison
\(\min\{1,t_k^{-1}\}\le\rho(t_k,d_k)\) forces \(t_k\to\infty\).
Also, whenever \(\rho(t_k,d_k)<1\),
\[
g_k^2\le\rho(t_k,d_k).
\]
In the nondense branch this follows from
\[
\log(1+d_k^2/t_k)
\le\log(e+L_kd_k^2/t_k)
\]
and removal of the cap when the rate is below one.
In the dense branch, using \(\log(1+x)\le x\) and
\(d_k^2/t_k\le1/L_k\) gives the same inequality.
Thus \(g_k\to0\).

Conversely, for every \(d\ge2\) and \(t>0\),
\[
\rho(t,d)\le
\left[
\frac{\log(1+d^2/t)}{L}
+\frac{\log(e+L)}{L}
\right]^2+\frac1{L^2}.
\]
For the nondense branch, this follows from
\[
e+L(d^2/t)\le(1+d^2/t)(e+L);
\]
for the dense branch, \(\rho(t,d)\le L^{-2}\).
Both correction terms tend to zero as \(d\to\infty\).

This bound also handles sequences with oscillating dimensions.
Given any positive tolerance, choose a dimension threshold so that both corrections are sufficiently small above that threshold.
Then \(g_k\to0\) makes the displayed upper bound small for all sufficiently large \(k\) whose dimensions exceed the threshold.
Below it, the elementary upper comparison gives
\[
\rho(t_k,d_k)\le\frac{4d_k^2}{t_k},
\]
which tends uniformly to zero as \(t_k\to\infty\).
Hence
\[
\rho(t_k,d_k)\to0
\quad\Longleftrightarrow\quad
t_k\to\infty\ \text{and}\ g_k\to0.
\]
% lean: CausalSmith.Stat.LdpOptvalueUniformFrontier.rho_consistency_iff

\item \textit{Parametric comparison and the quadratic-resource regime.}
Suppose that fixed positive finite constants
\(c_{\mathrm{cmp}}\le C_{\mathrm{cmp}}\) satisfy, eventually,
\[
\frac{c_{\mathrm{cmp}}}{t_k}
\le\rho(t_k,d_k)
\le\frac{C_{\mathrm{cmp}}}{t_k}.
\]
Since \(\rho(t_k,d_k)\le1\) and
\(t_k\rho(t_k,d_k)\ge\min\{t_k,d_k\}/4\), these imply
\[
t_k\ge c_{\mathrm{cmp}},
\qquad
\min\{t_k,d_k\}\le4C_{\mathrm{cmp}}
\]
eventually.

Conversely, suppose that there are \(a_{\mathrm{par}}>0\) and
\(M_{\mathrm{par}}\in\mathbb R\) such that, eventually,
\[
t_k\ge a_{\mathrm{par}},
\qquad
\min\{t_k,d_k\}\le M_{\mathrm{par}}.
\]
Define
\[
B_{\mathrm{par}}=\max\{1,M_{\mathrm{par}}\}.
\]
The elementary lower comparison gives
\[
\frac{\min\{1,a_{\mathrm{par}}\}}{t_k}
\le\min\{1,t_k^{-1}\}\le\rho(t_k,d_k).
\]
For the upper comparison, split according to \(t_k\le d_k\) or \(d_k\le t_k\).
In the first case \(t_k\le B_{\mathrm{par}}\), so
\(\rho(t_k,d_k)\le1\le4B_{\mathrm{par}}^2/t_k\).
In the second case \(d_k\le B_{\mathrm{par}}\), and
\[
\rho(t_k,d_k)\le\frac{4d_k^2}{t_k}
\le\frac{4B_{\mathrm{par}}^2}{t_k}.
\]
This proves the asserted general parametric equivalence.
When \(t_k\to\infty\), the positive lower-resource condition holds eventually. Moreover,
\[
\min\{t_k,d_k\}\le M_{\mathrm{par}}
\]
eventually forces \(d_k\le M_{\mathrm{par}}\) eventually, since
\(t_k>M_{\mathrm{par}}\) eventually. Enlarging to an integer bound proves the bounded-dimension characterization; its converse follows from
\(\min\{t_k,d_k\}\le d_k\).

For the quadratic-resource regime, fix
\(0<a_{\mathrm{el}}\le b_{\mathrm{el}}\).
Choose an integer \(D_0\) large enough that, whenever \(d\ge D_0\),
\[
L\ge b_{\mathrm{el}}+1,\qquad
\log L\ge2\log b_{\mathrm{el}},\qquad
\log L\ge\log(e+a_{\mathrm{el}}^{-1}).
\]
For every \(t\) in the specified band, define
\[
u_{\mathrm{el}}=\frac{t}{d^2}\in[a_{\mathrm{el}},b_{\mathrm{el}}].
\]
Then \(t<d^2L\), and
\[
\log(e+L/u_{\mathrm{el}})
\ge\log(L/b_{\mathrm{el}})
\ge\frac12\log L,
\]
\[
\log(e+L/u_{\mathrm{el}})
\le\log\!\left(L(e+a_{\mathrm{el}}^{-1})\right)
\le2\log L.
\]
Since \(0\le(\log L)/L\le1\), applying the cap and squaring gives
\[
\frac14\left(\frac{\log L}{L}\right)^2
\le\rho(t,d)
\le4\left(\frac{\log L}{L}\right)^2.
\]
These bounds supply the required constants uniformly over the band.
% lean: CausalSmith.Stat.LdpOptvalueUniformFrontier.rate_resource_conclusions

\item \textit{Transfer to the actual minimax values.}
The upper bounds already proved make the minimax values finite, so their real values satisfy the same four inequalities. For a fixed class and an admissible sequence, define
\[
\begin{aligned}
\mathcal R_{\mathcal C,k}^{\mathrm{TV}}
&=\mathfrak R^{\mathrm{TV}}_{\mathcal C}
(\mathcal P_{d_k}^{\mathrm{pair}};n_k,d_k,\varepsilon_k),\\
\mathcal H_{\mathcal C,k}^{\mathrm{TV}}
&=\mathfrak H^{\mathrm{TV}}_{\mathcal C}
(\mathcal P_{d_k}^{\mathrm{pair}};n_k,d_k,\varepsilon_k).
\end{aligned}
\]
The common constants give
\[
c\rho(t_k,d_k)\le\mathcal R_{\mathcal C,k}^{\mathrm{TV}}
\le C_0\rho(t_k,d_k),
\]
\[
c\sqrt{\rho(t_k,d_k)}
\le\mathcal H_{\mathcal C,k}^{\mathrm{TV}}
\le C_0\sqrt{\rho(t_k,d_k)}.
\]
Squeezing nonnegative sequences therefore yields
\[
\mathcal R_{\mathcal C,k}^{\mathrm{TV}}\to0
\Longleftrightarrow\rho(t_k,d_k)\to0,
\]
\[
\mathcal H_{\mathcal C,k}^{\mathrm{TV}}\to0
\Longleftrightarrow\sqrt{\rho(t_k,d_k)}\to0
\Longleftrightarrow\rho(t_k,d_k)\to0.
\]
Combining these equivalences with the rate consistency criterion proves both minimax consistency assertions.

Two-sided comparison is symmetric and transitive: reversing
\(c\rho\le\mathcal R\le C_0\rho\) gives
\(\mathcal R/C_0\le\rho\le\mathcal R/c\).
Consequently,
\[
\mathcal R_{\mathcal C,k}^{\mathrm{TV}}\asymp t_k^{-1}
\Longleftrightarrow
\rho(t_k,d_k)\asymp t_k^{-1}.
\]
The general and diverging-resource characterizations thus transfer directly.

Finally, in the band from the preceding step, the common bounds imply
\[
\frac c4\left(\frac{\log L}{L}\right)^2
\le
\mathfrak R^{\mathrm{TV}}_{\mathcal C}
(\mathcal P_d^{\mathrm{pair}};n,d,\varepsilon)
\le
4C_0\left(\frac{\log L}{L}\right)^2,
\]
and
\[
\frac c2\left|\frac{\log L}{L}\right|
\le
\mathfrak H^{\mathrm{TV}}_{\mathcal C}
(\mathcal P_d^{\mathrm{pair}};n,d,\varepsilon)
\le
2C_0\left|\frac{\log L}{L}\right|.
\]
To give the risk and length bounds the coupled constants in the statement, define
\[
c_{\mathrm{el}}=\min\{c/4,c^2/4\},
\qquad
C_{\mathrm{el}}=\max\{4C_0,4C_0^2\}.
\]
Then \(0<c_{\mathrm{el}}\le C_{\mathrm{el}}<\infty\), and
\[
\sqrt{c_{\mathrm{el}}}\le c/2,
\qquad
2C_0\le\sqrt{C_{\mathrm{el}}}.
\]
The displayed risk and length inequalities therefore imply exactly the final two bounds, with \(c_{\mathrm{el}},C_{\mathrm{el}}\) as the band-dependent constants. This completes all the sequence conclusions.
% lean: CausalSmith.Stat.LdpOptvalueUniformFrontier.value_sequence_conclusions_of_bounds
\end{enumerate}
\end{proof}
